% The Q arrow on a hatched cut face: three treatments compared at close-up.
\documentclass[10pt,border=1mm]{standalone}
\usepackage[T1]{fontenc}
\usepackage{figurestyle}
\input{primitives.tex}
\begin{document}
\begin{tikzpicture}[plate]
\path[use as bounding box] (0,0) rectangle (12,4.2);
\newcommand\face[2]{\begin{scope}[shift={(#1,2.3)}]
  \path[cut] (-.55,-1.2) rectangle (.55,1.2); \draw[heavy] (-.55,-1.2) rectangle (.55,1.2);
  \path[fill=ColBand] (-.08,-1.2) rectangle (.08,1.2); \draw[line width=.4pt,draw=PlateInk] (-.08,-1.2)--(-.08,1.2) (.08,-1.2)--(.08,1.2);
  \hatom{0}{1.2}{.07}{ColDot}
  #2
  \node[lab,anchor=south] at (0,1.32) {$\eta$};
  \end{scope}}
\face{1.6}{\draw[harrow] (0,.3)--(-.5,.3); \node[lab,anchor=east,inner sep=1pt] at (-.72,.3) {$Q$};}
\node[sub,anchor=north,align=center] at (1.6,.7) {A: arrow on the hatch,\\label outside, no box};
\face{5.0}{\draw[line width=3pt,draw=white] (0,.3)--(-.54,.3); \draw[harrow] (0,.3)--(-.5,.3); \node[lab,anchor=east,inner sep=1pt] at (-.72,.3) {$Q$};}
\node[sub,anchor=north,align=center] at (5.0,.7) {B: 3 pt white underlay};
\face{8.4}{\draw[line width=1.6pt,draw=white] (0,.3)--(-.5,.3); \draw[harrow] (0,.3)--(-.5,.3); \node[lab,anchor=east,inner sep=1pt] at (-.72,.3) {$Q$};}
\node[sub,anchor=north,align=center] at (8.4,.7) {C: 1.6 pt underlay,\\head on the hatch};
\end{tikzpicture}
\end{document}
